% GENERATED FILE. Edit canonical lessons, then run npm run generate:books.
% canonical-source-hash: fnv1a64:4c2aef73c24ca553
% canonical-lessons: GU-C158-saph-karvu, GU-C158-valvu, GU-C158-luchvu, GU-C158-potu, GU-C158-dagh, GU-R158-review-153-158

\chapter{Sweeping and wiping}
\label{ch:gu-sweeping-and-wiping}
\hlchaptermodality{158}

\begin{chapteropening}
\textbf{By the end of this chapter:} \emph{I can describe how we sweep, wipe and mop the house.}

Complete the last lesson of chapter 158: I can describe how we sweep, wipe and mop the house.
\end{chapteropening}

\section[\texorpdfstring{\gu{સાફ} \gu{કરવું} (sāph karvũ)}{સાફ કરવું (sāph karvũ)}]{\gu{સાફ} \gu{કરવું} (sāph karvũ) — to clean up, to clean}

\label{lesson:GU-C158-saph-karvu}



\begin{quote}
Before the new one: say the Gujarati for curtain, then the Gujarati for shelf.
\end{quote}

\subsection*{You'll want to know: \gu{સાફ} \gu{કરવું}}
\textbf{\gu{સાફ} \gu{કરવું}} — \emph{sāph karvũ} — \textquotedblleft{}to clean up, to clean\textquotedblright{}.

\begin{quote}
\textbf{\gu{અમે} \gu{રોજ} \gu{ઘર} \gu{સાફ} \gu{કરીએ} \gu{છીએ}.} — \emph{ame roj ghar sāph karīe chīe} — \textquotedblleft{}We clean the house every day.\textquotedblright{}
\end{quote}

\begin{cousinweb}
\textbf{\gu{સાફ} \gu{કરવું}} (\emph{sāph karvũ}) is built from \emph{sāph}, \textquotedblleft{}clean\textquotedblright{}, and the inherited \emph{karvũ}, \textquotedblleft{}to do\textquotedblright{}. \emph{Sāph} is a loan from Persian \emph{sāf}, \textquotedblleft{}clear, pure\textquotedblright{} (EtymDB), which comes from Arabic; its Arabic root is \emph{ṣ-f-w}, \textquotedblleft{}to be clear\textquotedblright{} (Lane).
\end{cousinweb}

\begin{grammarlens}[title={What you've built}]
One more word for everyday life at A2.
\end{grammarlens}

\subsection*{Your turn}
\begin{itemize}
\raggedright
  \item \emph{Say it:} \emph{sāph karvũ}
  \item \emph{Say it:} \emph{sāph karvũ}, once more
  \item \emph{Say it:} \emph{ame roj ghar sāph karīe chīe} — We clean the house every day.
  \item \emph{Recall:} say the Gujarati for curtain, then the Gujarati for shelf, then say \emph{sāph karvũ} again
\end{itemize}

\begin{tcolorbox}[breakable,colback=teal!4,colframe=teal!35!black,title={Before you move on}]
What does \gu{સાફ} \gu{કરવું} mean? (\textquotedblleft{}To clean up, to clean\textquotedblright{}.) Say it once more.
\end{tcolorbox}
\section[\texorpdfstring{\gu{વાળવું} (vāḷvũ)}{વાળવું (vāḷvũ)}]{\gu{વાળવું} (vāḷvũ) — to sweep, to fold}

\label{lesson:GU-C158-valvu}



\begin{quote}
Before the new one: say the Gujarati for shelf, then the Gujarati for to clean up.
\end{quote}

\subsection*{You'll want to know: \gu{વાળવું}}
\textbf{\gu{વાળવું}} — \emph{vāḷvũ} — \textquotedblleft{}to sweep, to fold\textquotedblright{}.

\begin{quote}
\textbf{\gu{તે} \gu{સાવરણીથી} \gu{કચરો} \gu{વાળે} \gu{છે}.} — \emph{te sāvarṇīthī kacro vāḷe che} — \textquotedblleft{}She sweeps the rubbish with a broom.\textquotedblright{}
\end{quote}

It also means to fold or to turn.

\begin{cousinweb}
\textbf{\gu{વાળવું}} (\emph{vāḷvũ}) is inherited, but the sweeping sense is not settled. CDIAL lists Gujarati \emph{vāḷvũ}, \textquotedblleft{}to fold, twist, turn\textquotedblright{}, under Sanskrit \emph{valati}, \textquotedblleft{}turns\textquotedblright{}, and does not list the sense of sweeping (CDIAL 11405).
\end{cousinweb}

\begin{grammarlens}[title={What you've built}]
One more word for everyday life at A2.
\end{grammarlens}

\subsection*{Your turn}
\begin{itemize}
\raggedright
  \item \emph{Say it:} \emph{vāḷvũ}
  \item \emph{Say it:} \emph{vāḷvũ}, once more
  \item \emph{Say it:} \emph{te sāvarṇīthī kacro vāḷe che} — She sweeps the rubbish with a broom.
  \item \emph{Recall:} say the Gujarati for shelf, then the Gujarati for to clean up, then say \emph{vāḷvũ} again
\end{itemize}

\begin{tcolorbox}[breakable,colback=teal!4,colframe=teal!35!black,title={Before you move on}]
What does \gu{વાળવું} mean? (\textquotedblleft{}To sweep, to fold\textquotedblright{}.) Say it once more.
\end{tcolorbox}
\section[\texorpdfstring{\gu{લૂછવું} (lūchvũ)}{લૂછવું (lūchvũ)}]{\gu{લૂછવું} (lūchvũ) — to wipe, to dry off}

\label{lesson:GU-C158-luchvu}



\begin{quote}
Before the new one: say the Gujarati for to clean up, then the Gujarati for to sweep.
\end{quote}

\subsection*{You'll want to know: \gu{લૂછવું}}
\textbf{\gu{લૂછવું}} — \emph{lūchvũ} — \textquotedblleft{}to wipe, to dry off\textquotedblright{}.

\begin{quote}
\textbf{\gu{હું} \gu{ટુવાલથી} \gu{હાથ} \gu{લૂછું} \gu{છું}.} — \emph{hũ ṭuvālthī hāth lūchũ chũ} — \textquotedblleft{}I wipe my hands with a towel.\textquotedblright{}
\end{quote}

\begin{cousinweb}
The origin of \textbf{\gu{લૂછવું}} (\emph{lūchvũ}) is not settled. CDIAL lists no Gujarati verb of this form; it lists the similar \emph{pũchvũ}, \textquotedblleft{}to wipe\textquotedblright{}, under Sanskrit \emph{proñchati}, \textquotedblleft{}wipes off\textquotedblright{}, without linking the two (CDIAL 9011).
\end{cousinweb}

\begin{grammarlens}[title={What you've built}]
One more word for everyday life at A2.
\end{grammarlens}

\subsection*{Your turn}
\begin{itemize}
\raggedright
  \item \emph{Say it:} \emph{lūchvũ}
  \item \emph{Say it:} \emph{lūchvũ}, once more
  \item \emph{Say it:} \emph{hũ ṭuvālthī hāth lūchũ chũ} — I wipe my hands with a towel.
  \item \emph{Recall:} say the Gujarati for to clean up, then the Gujarati for to sweep, then say \emph{lūchvũ} again
\end{itemize}

\begin{tcolorbox}[breakable,colback=teal!4,colframe=teal!35!black,title={Before you move on}]
What does \gu{લૂછવું} mean? (\textquotedblleft{}To wipe, to dry off\textquotedblright{}.) Say it once more.
\end{tcolorbox}
\section[\texorpdfstring{\gu{પોતું} (potũ)}{પોતું (potũ)}]{\gu{પોતું} (potũ) — floor cloth, mop}

\label{lesson:GU-C158-potu}



\begin{quote}
Before the new one: say the Gujarati for to sweep, then the Gujarati for to wipe.
\end{quote}

\subsection*{You'll want to know: \gu{પોતું}}
\textbf{\gu{પોતું}} — \emph{potũ} — \textquotedblleft{}floor cloth, mop\textquotedblright{}.

\begin{quote}
\textbf{\gu{તે} \gu{સવારે} \gu{પોતું} \gu{કરે} \gu{છે}.} — \emph{te savāre potũ kare che} — \textquotedblleft{}He mops the floor in the morning.\textquotedblright{}
\end{quote}

\textbf{\gu{પોતું} \gu{કરવું}}, \emph{potũ karvũ}, to mop the floor.

\begin{cousinweb}
\textbf{\gu{પોતું}} (\emph{potũ}) is a tadbhava. It comes from \emph{*pōtta}, \textquotedblleft{}cloth\textquotedblright{}, a reconstructed by-form of Sanskrit \emph{pōta}, \textquotedblleft{}cloth\textquotedblright{}, and CDIAL glosses the Gujarati word \textquotedblleft{}rag\textquotedblright{} (CDIAL 8400).
\end{cousinweb}

\begin{grammarlens}[title={What you've built}]
One more word for everyday life at A2.
\end{grammarlens}

\subsection*{Your turn}
\begin{itemize}
\raggedright
  \item \emph{Say it:} \emph{potũ}
  \item \emph{Say it:} \emph{potũ}, once more
  \item \emph{Say it:} \emph{te savāre potũ kare che} — He mops the floor in the morning.
  \item \emph{Recall:} say the Gujarati for to sweep, then the Gujarati for to wipe, then say \emph{potũ} again
\end{itemize}

\begin{tcolorbox}[breakable,colback=teal!4,colframe=teal!35!black,title={Before you move on}]
What does \gu{પોતું} mean? (\textquotedblleft{}Floor cloth, mop\textquotedblright{}.) Say it once more.
\end{tcolorbox}
\section[\texorpdfstring{\gu{ડાઘ} (ḍāgh)}{ડાઘ (ḍāgh)}]{\gu{ડાઘ} (ḍāgh) — stain, spot}

\label{lesson:GU-C158-dagh}



\begin{quote}
Before the new one: say the Gujarati for to wipe, then the Gujarati for floor cloth.
\end{quote}

\subsection*{You'll want to know: \gu{ડાઘ}}
\textbf{\gu{ડાઘ}} — \emph{ḍāgh} — \textquotedblleft{}stain, spot\textquotedblright{}.

\begin{quote}
\textbf{\gu{આ} \gu{ડાઘ} \gu{નથી} \gu{જતો}.} — \emph{ā ḍāgh nathī jato} — \textquotedblleft{}This stain does not come out.\textquotedblright{}
\end{quote}

\begin{cousinweb}
\textbf{\gu{ડાઘ}} (\emph{ḍāgh}) is a loan, usually traced to Persian \emph{dāgh}, \textquotedblleft{}mark, brand, stain\textquotedblright{}. CDIAL, which follows inherited words, lists no Gujarati form.
\end{cousinweb}

\begin{grammarlens}[title={What you've built}]
One more word for everyday life at A2.
\end{grammarlens}

\subsection*{Your turn}
\begin{itemize}
\raggedright
  \item \emph{Say it:} \emph{ḍāgh}
  \item \emph{Say it:} \emph{ḍāgh}, once more
  \item \emph{Say it:} \emph{ā ḍāgh nathī jato} — This stain does not come out.
  \item \emph{Recall:} say the Gujarati for to wipe, then the Gujarati for floor cloth, then say \emph{ḍāgh} again
\end{itemize}

\begin{tcolorbox}[breakable,colback=teal!4,colframe=teal!35!black,title={Before you move on}]
What does \gu{ડાઘ} mean? (\textquotedblleft{}Stain, spot\textquotedblright{}.) Say it once more.
\end{tcolorbox}
\section[(dialogue)]{Home and dress words, a review}

\label{lesson:GU-R158-review-153-158}



\begin{quote}
Say the Gujarati for room and kitchen.
\end{quote}

\subsection*{Your turn}
Say these aloud:

\begin{itemize}
\raggedright
  \item room, kitchen, terrace, veranda, stairs
  \item cloth, sari, scarf, to take off, to wrap oneself in
  \item to wake someone, to get ready, to comb one's hair, perfume, to apply
  \item ornament, bindi, henna, to decorate, to tie
  \item to arrange, to put, to hang up, curtain, shelf
  \item to clean up, to sweep, to wipe, floor cloth, stain
\end{itemize}

\begin{tcolorbox}[breakable,colback=teal!4,colframe=teal!35!black,title={Before you move on}]
Room? (\textbf{\gu{ઓરડો}}.) Kitchen? (\textbf{\gu{રસોડું}}.) Terrace? (\textbf{\gu{અગાસી}}.) Veranda? (\textbf{\gu{ઓસરી}}.) Stairs? (\textbf{\gu{દાદર}}.) Cloth? (\textbf{\gu{કપડું}}.) Sari? (\textbf{\gu{સાડી}}.) Scarf? (\textbf{\gu{દુપટ્ટો}}.) To take off? (\textbf{\gu{ઉતારવું}}.) To wrap oneself in? (\textbf{\gu{ઓઢવું}}.) To wake someone? (\textbf{\gu{ઉઠાડવું}}.) To get ready? (\textbf{\gu{તૈયાર} \gu{થવું}}.) To comb one's hair? (\textbf{\gu{ઓળવું}}.) Perfume? (\textbf{\gu{અત્તર}}.) To apply? (\textbf{\gu{લગાડવું}}.) Ornament? (\textbf{\gu{ઘરેણું}}.) Bindi? (\textbf{\gu{ચાંલ્લો}}.) Henna? (\textbf{\gu{મેંદી}}.) To decorate? (\textbf{\gu{સજાવવું}}.) To tie? (\textbf{\gu{બાંધવું}}.) To arrange? (\textbf{\gu{ગોઠવવું}}.) To put? (\textbf{\gu{મૂકવું}}.) To hang up? (\textbf{\gu{ટાંગવું}}.) Curtain? (\textbf{\gu{પડદો}}.) Shelf? (\textbf{\gu{છાજલી}}.) To clean up? (\textbf{\gu{સાફ} \gu{કરવું}}.) To sweep? (\textbf{\gu{વાળવું}}.) To wipe? (\textbf{\gu{લૂછવું}}.) Floor cloth? (\textbf{\gu{પોતું}}.) Stain? (\textbf{\gu{ડાઘ}}.)
\end{tcolorbox}
